\section{Các Công Trình Liên Quan}
\label{sec:related}

Phần này thảo luận về các phương pháp tiên tiến được phân loại theo lọc, phát hiện bất thường,
theo dõi xuất xứ và tăng cường bảo mật kernel, làm nổi bật điểm mạnh và hạn chế của chúng so
với khung DeSFAM. Bảng~\ref{tab:related_work} tóm tắt các phương pháp lọc lời gọi hệ thống và
phát hiện bất thường gần đây, làm nổi bật trọng tâm vấn đề, hạn chế và kỹ thuật cốt lõi của từng phương pháp.

\begin{table*}[tp]
  \centering
  \caption{So sánh các giải pháp bảo mật container tiên tiến.}
  \label{tab:related_work}
  \includegraphics[width=\linewidth]{table1}
\end{table*}

\subsection{Các Phương Pháp Lọc Lời Gọi Hệ Thống}
\label{subsec:syscall_filtering}

Confine~\cite{ref31_confine} sử dụng phân tích tĩnh để xây dựng các hồ sơ seccomp hạn chế,
đạt tỷ lệ giảm lời gọi hệ thống (SRR) là 47,6\%. Tuy nhiên, tính chất tĩnh của nó ngăn cản
sự thích nghi với các hành vi thời gian chạy, làm giảm hiệu quả trước các khối lượng công việc
luôn thay đổi. Tương tự, Prof-gen \cite{ref36_profgen} dựa vào các kỹ thuật tĩnh đòi hỏi cấu
hình thủ công hoặc bỏ qua theo dõi động, hạn chế khả năng mở rộng và độ bền trong các triển
khai thực tế.

KubeRosy~\cite{ref29_kuberosy} và Optimus~\cite{ref27_optimus} tích hợp eBPF để giám sát lời
gọi hệ thống trực tiếp và cập nhật chính sách thời gian chạy. KubeRosy tận dụng Kubernetes CRD
để tinh chỉnh chính sách theo thời gian thực nhưng vẫn bị ràng buộc với môi trường Kubernetes.
Optimus sử dụng lọc dựa trên liên kết để thích nghi lặp đi lặp lại các hồ sơ lời gọi hệ thống,
tuy nhiên tính tự động của nó bị hạn chế---yêu cầu sự phê duyệt của người vận hành đối với các
lời gọi hệ thống không khớp và dừng cập nhật hồ sơ nếu phân tích liên kết thất bại. Prof-gen,
mặc dù kết hợp theo dõi động lúc khởi động với phân tích đồ thị lời gọi tĩnh để tạo danh sách
truy cập tối thiểu mà không cần mã nguồn, vẫn thiếu khả năng nắm bắt các thay đổi hành vi trong
quá trình thực thi thời gian chạy.

\subsection{Hệ Thống Phát Hiện Xâm Nhập Dựa Trên Bất Thường}
\label{subsec:anomaly_ids}

Các kỹ thuật dựa trên học máy đã được nghiên cứu rộng rãi để phát hiện bất thường lời gọi hệ
thống. IDS dựa trên BoSC sử dụng mô hình hóa Túi Lời Gọi Hệ Thống với các bộ phân loại như
Rừng Ngẫu Nhiên, mang lại độ chính xác hợp lý trong môi trường quy mô nhỏ. Tuy nhiên, yêu cầu
huấn luyện lại thường xuyên cản trở khả năng mở rộng và thích nghi.

Các mô hình nhẹ như Bộ Phát Hiện Thoát dựa trên TinyML nhắm đến các thiết bị IoT và biên nhưng
không đủ khả năng xử lý các véc-tơ tấn công phức tạp trong các hệ thống đám mây hiệu suất cao
do tính thích nghi và độ chính xác hạn chế.

DeSFAM giải quyết những hạn chế này trong Giai đoạn~2 thông qua bộ phát hiện bất thường lai kết
hợp Bộ Mã Hóa Tự Động Biến Phân (VAE) với Rừng Cô Lập, được huấn luyện trên tập dữ liệu
DongTing~\cite{ref26_dongtingdataset}. Phương pháp này cho phép phát hiện chính xác, độ trễ thấp
đối với các cuộc tấn công zero-day mà không cần huấn luyện lại, đồng thời hỗ trợ suy luận cửa
sổ trượt thời gian thực.

\subsection{Các Phương Pháp Phát Hiện Dựa Trên Xuất Xứ}
\label{subsec:provenance}

Các hệ thống hướng xuất xứ như PACED giám sát các mối quan hệ lời gọi hệ thống và quan hệ nhân
quả giữa các container để phát hiện các cuộc tấn công xuyên không gian tên. Mặc dù hiệu quả,
chúng phải chịu chi phí cao do xây dựng đồ thị liên tục và khả năng mở rộng hạn chế.

DeSFAM vượt qua điều này bằng cách nhúng ngữ cảnh lời gọi hệ thống (ví dụ: PID, không gian tên,
dấu thời gian) trực tiếp vào các đặc trưng mô hình, cho phép nhận thức hành vi mà không phải
chịu chi phí của việc theo dõi xuất xứ đầy đủ.

\subsection{Các Cải Tiến Bảo Mật Cấp Kernel}
\label{subsec:kernel_security}

Các biện pháp phòng thủ tầng kernel như Condo sử dụng Môi Trường Thực Thi Đáng Tin Cậy (TEE) để
bảo vệ tính toàn vẹn quyền hạn. Mặc dù an toàn, sự phụ thuộc vào phần cứng chuyên dụng và sửa
đổi kernel khiến chúng không phù hợp cho các triển khai đám mây quy mô lớn.

Ngược lại, DeSFAM đạt được tính di động và khả năng mở rộng bằng cách sử dụng các bản đồ eBPF
và hook LSM để thực thi các chính sách lời gọi hệ thống một cách động, mà không cần sửa đổi
kernel hay phụ thuộc vào TEE.

\subsection{Các Hạn Chế Chung Của Các Phương Pháp Hiện Tại}
\label{subsec:common_limitations}

Mặc dù đã có những tiến bộ đáng kể trong việc tối thiểu hóa lời gọi hệ thống và bảo mật
container, các giải pháp hiện tại vẫn đối mặt với một số hạn chế chính:

\begin{itemize}
  \item Các chính sách tĩnh không thể thích nghi với các hành vi thời gian chạy luôn thay đổi.
  \item Thiếu thực thi thời gian thực cho phép lạm dụng lời gọi hệ thống giữa các lần cập nhật.
  \item Các phương pháp học máy yêu cầu huấn luyện lại thường xuyên, hạn chế khả năng mở rộng.
  \item Các phụ thuộc phần cứng (ví dụ: TEE) hạn chế tính linh hoạt trong triển khai.
\end{itemize}

\subsection{DeSFAM: Thúc Đẩy Trình Độ Hiện Tại}
\label{subsec:desfam_advancement}

DeSFAM giải quyết các hạn chế chính của các hệ thống trước thông qua một kiến trúc đa giai đoạn
thống nhất. Nó kết hợp lập hồ sơ lời gọi hệ thống lai, kết hợp phân tách nhị phân tĩnh với theo
dõi thời gian chạy động, để tạo ra các danh sách truy cập lời gọi hệ thống chính xác. Thực thi
thời gian thực được thực hiện thông qua các probe eBPF và hook LSM, cho phép chặn độ trễ thấp
và cập nhật chính sách mà không cần khởi động lại container. Phát hiện bất thường được xử lý
bởi SyscallAD, một mô hình lai sử dụng Bộ Mã Hóa Tự Động Biến Phân và Rừng Cô Lập để phát hiện
chính xác, không cần huấn luyện lại đối với các cuộc tấn công đã biết và zero-day. DeSFAM cũng
tích hợp lọc nhận thức CVE và chấm điểm rủi ro dựa trên MITRE ATT\&CK để hạn chế động các lời
gọi hệ thống dễ bị tổn thương, đồng thời áp dụng logic phản hồi phân tầng để thực thi các hành
động chi tiết dựa trên mức độ nghiêm trọng của mối đe dọa. Khung hoạt động mà không cần sửa đổi
kernel hay phụ thuộc vào các nền tảng điều phối, cung cấp bảo mật lời gọi hệ thống có khả năng
mở rộng và thích nghi cho các môi trường container hiện đại.
